\section{Preliminaries}
\label{sec:preliminaries}

This section fixes notation, recalls the definitions used throughout,
and isolates one linear-algebra fact that makes the certified system in
\cref{sec:certification} square.

\subsection{Cubic graphs and $\Sphere$-flows}
\label{sec:prelim-flows}

All graphs are finite, simple, and unoriented. A graph $G$ is
\emph{cubic} if every vertex has degree $3$, and \emph{bridgeless} if
removing any single edge leaves $G$ connected. For each edge
$e = \{u,v\}$ we fix one orientation once and for all; the sign
$\sigma(v,e) \in \{\pm 1\}$ records whether $v$ is the tail
($\sigma = +1$) or the head ($\sigma = -1$) of the oriented edge. The
codebase fixes the orientation deterministically by the canonical edge
order: for $e = \{u,v\}$ with $u < v$ in the integer labelling,
$\sigma(u, e) = +1$ and $\sigma(v, e) = -1$
(\path{scripts/witness.py:orient}).

For an abelian group $A$, an \emph{$A$-flow} on $G$ is a map
$\varphi : E(G) \to A$ satisfying Kirchhoff's law
$\sum_{e \ni v} \sigma(v,e)\,\varphi(e) = 0$ at every vertex
$v \in V(G)$. A flow is \emph{nowhere-zero} if $\varphi(e) \ne 0$ for
every $e$. Replacing $A$ with the sphere
$\Sphere \subset \R^{3}$ (taking values in a set rather than a group;
the Kirchhoff sum is interpreted in $\R^{3}$) gives the central object
of this paper:

\begin{definition}[$\Sphere$-flow]
\label{def:s2-flow}
An \emph{$\Sphere$-flow} on $G$ is a map $\varphi : E(G) \to \Sphere$
such that
\[
   \sum_{e \ni v} \sigma(v,e)\,\varphi(e) \;=\; 0
   \qquad \text{for every } v \in V(G),
\]
where the sum is taken in $\R^{3}$. Equivalently (cf.\
\cref{eq:kirchhoff}), $\varphi$ assigns a unit vector in $\R^{3}$ to
each oriented edge so that Kirchhoff's law holds at every vertex.
The nowhere-zero condition is automatic since unit vectors are nonzero.
\end{definition}

\begin{lemma}[120-degree triple, after~{\cite[Observation~9]{HMM26}}]
\label{lem:120-degree}
For unit vectors $v_1, v_2, v_3 \in \R^{3}$,
$v_1 + v_2 + v_3 = 0$ if and only if the four-point set
$\{0, v_1, v_2, v_3\}$ is coplanar and every pair of the $v_i$ makes
angle $2\pi/3$.
\end{lemma}

\begin{proof}
($\Leftarrow$) If the three unit vectors are coplanar and pairwise at
angle $2\pi/3$, then in the common plane they form the three cube
roots of unity (up to a rotation), whose sum is zero.
($\Rightarrow$) From $v_1 + v_2 + v_3 = 0$, squaring gives
$3 + 2(v_1{\cdot}v_2 + v_1{\cdot}v_3 + v_2{\cdot}v_3) = 0$, so the
three pairwise inner products sum to $-3/2$. Each is at most $1$ in
absolute value, so each must equal $-1/2$, i.e.\ each pair makes
angle $2\pi/3$. Coplanarity follows from $v_3 = -v_1 - v_2$.
\end{proof}

In particular, for a cubic graph the Kirchhoff equation at every
vertex specialises to the constraint that the three incident unit
vectors lie in a common plane through the origin at pairwise
$2\pi/3$.

\subsection{Snarks and nontrivial snarks}
\label{sec:prelim-snarks}

The \emph{cyclic edge-connectivity} $\lambda_c(G)$ of a cubic graph
$G$ with at least two disjoint cycles is the minimum size of an edge
cut both of whose sides contain a cycle. A \emph{snark} is a simple
connected cubic bridgeless graph with chromatic index $\chi'(G) = 4$
(equivalently, by Vizing's theorem on cubic graphs, $G$ is not
$3$-edge-colourable).

\begin{definition}[Nontrivial snark]
\label{def:nontrivial-snark}
A snark $G$ is \emph{nontrivial} if it has girth at least $5$ and is
cyclically $4$-edge-connected ($\lambda_c(G) \ge 4$).
\end{definition}

This is the convention of Brinkmann, Goedgebeur, H{\"a}gglund, and
Markstr{\"o}m~\cite{BGHM13}. The catalogue used in \cref{sec:finite-theorem}
is exactly the enumeration of nontrivial snarks on at most $28$
vertices from~\cite{BGHM13} (with the $n = 28$ count produced by the
\texttt{snarkhunter}~\cite{snarkhunter} generator of
Brinkmann--Goedgebeur--McKay).

\subsection{Cycle double covers and $H_d$-flows}
\label{sec:prelim-cdc}

A \emph{cycle double cover (CDC)} of a graph $G$ is a multiset of
cycles such that every edge is contained in exactly two of them.
An \emph{orientation} of a CDC assigns a direction to each cycle so
that every edge is traversed once in each direction by the two cycles
through it. For an integer $d \ge 2$, an \emph{$H_d$-flow} on a cubic
graph (in the sense of Mattiolo, Mazzuoccolo, Rajnik, and
Tabarelli~\cite{MMRT25}) is an assignment of edge vectors taken from
the boundary group of the $(d-1)$-cube. For $d = 4$,
\cite{MMRT25}~Theorem~4 states that on cubic graphs an oriented
$4$-CDC is equivalent to an $H_4$-flow, which restricted to coordinates
modulo $2$ gives a nowhere-zero $\Z_2^{2}$-flow, hence a nowhere-zero
$4$-flow. We use this only as a sharp negative calibration in
\cref{sec:cdc-obstruction}: a snark, having $\chi'(G) = 4$, admits no
nowhere-zero $4$-flow and therefore no oriented $4$-CDC.

\TODO{Calibration note: the description above paraphrases MMRT25
Theorem 4. The exact statement and proof structure are not yet
transcribed from the source paper; see \path{docs/literature_notes.md}
and \path{docs/cdc_negative_calibration.md} for the project's own
notes.}

\subsection{The pinning gauge}
\label{sec:prelim-gauge}

The $\Sphere$-flow equations are equivariant under the diagonal
action of $\SO(3)$: if $\varphi$ is an $\Sphere$-flow then so is
$R \circ \varphi$ for any $R \in \SO(3)$. To obtain a system with
finitely many real solutions (and hence to apply the Krawczyk
operator), we fix this $3$-parameter symmetry by \emph{pinning} the
three spokes at vertex $0$ to the canonical $120$-degree triple in the
$xy$-plane,
\[
   \varphi(e_1) \;=\; (1, 0, 0), \qquad
   \varphi(e_2) \;=\; (-\tfrac{1}{2},\, \tfrac{s}{2},\, 0), \qquad
   \varphi(e_3) \;=\; (-\tfrac{1}{2},\, -\tfrac{s}{2},\, 0),
\]
where $e_1, e_2, e_3$ are the three edges incident to vertex $0$ in
canonical order and the formal symbol $s$ is constrained by the
auxiliary relation $s^{2} - 3 = 0$ so that $s$ represents $\sqrt{3}$
exactly throughout the rational computation. See
\path{scripts/exact.py:build_ideal} for the implementation.

The pinning has two simultaneous consequences. First, it eliminates
nine of the original $3|E|$ scalar variables (the three
coordinates of each of the three pinned edges). Second, the Kirchhoff
equation at vertex $0$ is now satisfied identically (the canonical
triple sums to zero), so it carries no further constraint.

The remaining gauge freedom is removed by dropping the Kirchhoff
equation at one further \emph{redundant} vertex (the implementation
drops the last vertex in canonical order). This is justified by
\cref{lem:square-system} below.

\begin{lemma}[Square pinned system]
\label{lem:square-system}
Let $G$ be a connected cubic graph with $|V|$ vertices and $|E| = 3|V|/2$
edges, and fix the orientation convention above. Then:
\begin{enumerate}[label=(\roman*)]
  \item The $|V|$ vector-valued vertex Kirchhoff equations carry
        exactly $3$ linear dependencies: the coordinate-wise sum of
        all Kirchhoff equations vanishes identically (one dependency
        per coordinate axis of $\R^{3}$).
  \item After pinning vertex $0$ as in the canonical triple above
        and dropping the (vector) Kirchhoff equation at one further
        vertex, the resulting polynomial system together with the
        $|E|-3$ unit-norm equations $\lVert \varphi(e) \rVert^{2} - 1$
        on the unpinned edges and the auxiliary relation $s^{2} - 3$
        has exactly $|E| - 3 + 1$ polynomials (counting $s^{2} - 3$)
        in the same number of free variables.
\end{enumerate}
The resulting system is \emph{square} in this sense, has total degree
at most $2$, and is the system fed to the Krawczyk operator in
\cref{sec:certification-krawczyk}.
\end{lemma}

\begin{proof}
(i) The orientation convention assigns each edge $e = \{u, v\}$ a
sign $\sigma(u, e) = +1$ at one endpoint and $\sigma(v, e) = -1$ at
the other, so $\sigma(u, e) + \sigma(v, e) = 0$ for every edge. For
any fixed coordinate $j \in \{1, 2, 3\}$ of $\R^{3}$, the sum
$\sum_{v \in V} \big( \sum_{e \ni v} \sigma(v, e)\,\varphi(e) \big)_j$
counts each edge exactly twice with opposite signs and therefore
vanishes identically as a polynomial in the edge variables. This
gives three independent linear dependencies among the $|V|$ vector
Kirchhoff equations, one per coordinate. No further dependency exists
on a connected cubic graph: the
incidence matrix of an oriented connected graph has rank $|V| - 1$
over the rationals, which combined with the three-coordinate
structure produces rank $3(|V| - 1) = 3|V| - 3$ in the linearised
Kirchhoff system.

(ii) Pinning eliminates the three vector variables on the three
spokes at vertex $0$ ($9$ scalar variables), and the Kirchhoff
equation at vertex $0$ is then identically satisfied. The remaining
$|V| - 1$ vector Kirchhoff equations contribute $3(|V| - 1)$ scalar
equations modulo the $3$ dependencies of part~(i). Two of those
dependencies cease to be free once vertex $0$ is pinned (the
canonical triple satisfies them by construction), so exactly one
linear dependency remains; dropping the Kirchhoff equation at any
single further vertex kills it. The result is $3(|V| - 1) - 3 =
3|V| - 6$ Kirchhoff scalar equations plus $|E| - 3$ unit-norm
equations plus the auxiliary $s^{2} - 3$, for a total of
$3|V| - 6 + |E| - 3 + 1 = 3|V| + |E| - 8$ polynomials. The free
variables are the $3(|E| - 3) = 3|E| - 9$ unpinned edge coordinates
together with $s$, for a total of $3|E| - 8$. Using $|E| = 3|V|/2$
both counts simplify to $9|V|/2 - 8$, confirming squareness.
\end{proof}

The bookkeeping in (ii) is the linear-algebra fact that makes the
square-system construction in \path{scripts/exact.py:build_ideal}
correct; in particular it is what justifies the call
\texttt{build\_ideal(G, pin=True, drop\_redundant=True)} used by the
certifier.
